\subsection{Shorter endpoint parameters and a mixed-type period construction}

\begin{proposition}[DP52: the exact endpoint and period interface]\label{DP52}
The following statements hold for the recovered source interfaces.
\begin{enumerate}[label=(\alph*),leftmargin=*]
\item Let $D,D'$ have the same $k$ and identified complex places, with even lengths
$2\le M'\le M$. Rebuild the smaller rank record and spaces as in DP44.
Put $r=k+M-1$, $r'=k+M'-1$.
For $P:\operatorname{TPar}(D)$, retain its $t,u,v,\eta$ and restrict $b,e$ to
their first $M'$ entries. This defines $P':\operatorname{TPar}(D')$.
Interlacing, parity and centredness pass to $P'$, and for $j\le M'$,
\[
 X_{P'}(j,\tau)=X_P(j,\tau),\qquad
 D_{P'}(j,\tau)\subseteq D_P(j,\tau),
\]
including the Boolean signs in $D_P$.
The noncompact-gap condition passes as well.
For any source theta packages $X,X'$ on these records,
\[
 \operatorname{Hyp3}(X,P),\qquad 1801\le k+M'-1+t
 \quad\Longrightarrow\quad \operatorname{Hyp3}(X',P').
\]
The remaining inequality follows, for example, from
$r+t\ge1801+(M-M')$.
\item The six statements forming $\operatorname{AdmS}(D,X)$, together with
$\operatorname{Hyp3}(X,P)$, imply the source's $\operatorname{PointsLemma}$
at $P$. Their finite downward induction works for every permitted even $M$.
Thus a rebuilt shorter family satisfying these same six statements has the
same nonzero-period conclusion; no additional minimum induction length is
needed.
\item The actual endpoint period is uniform in the rank record.
More strongly, its row and point goodness certificates and its inner-product
formula can be constructed without $\operatorname{Fields.Ok.quad}$.
It suffices to assume the exact conjugation compatibilities
\[
 \phi(\bar x)=\overline{\phi(x)}
       \quad(\phi:E\hookrightarrow\mathbf C),\qquad
 \sigma_0(\bar y)=\overline{\sigma_0(y)}
       \quad(y\in K),
\]
for one coefficient-field embedding $\sigma_0$, and
\[
 \operatorname{emb}_\tau(\bar x)
       =\overline{\operatorname{emb}_\tau(x)},
\]
together with the specified weight and endpoint-level inputs.
These assumptions hold for $E=K=\mathbf Q(\zeta_7)$,
$F=E^+$ and the actual type $\{1,2,3\}$.
Consequently that mixed type has the actual finite-class row/point spaces,
endpoint period and inner-product formula at every admissible rebuilt rank.
This assertion does not say that the period is nonzero on a prescribed
iterated theta image.
\item The unchanged field wrapper nevertheless fails for that type:
\[
 E_{\{1,2,3\}}
 :=\{z\in E:\sigma_1(z)=\sigma_2(z)=\sigma_3(z)\}
 =\mathbf Q,
 \qquad \sigma_a(\zeta_7)=\zeta_7^a.
\]
Thus its field $\operatorname{Fields.Kq}$ has degree one over $\mathbf Q$,
whereas $\operatorname{Fields.Ok.quad}$ requires degree two.
This failure is independent of $M$ and persists under coefficient-field
enlargement, simultaneous relabeling and conjugation of the type.
For a Galois CM field, the degree-two condition is equivalent to the type
being induced from one embedding of an imaginary quadratic subfield.
\end{enumerate}
Part (c) explicitly bypasses the failed wrapper at the endpoint period-map
layer. It does not construct the full mixed-type theta correspondence or
its selected quantitative certificates.
\end{proposition}

\begin{proof}
\textit{Parameter restriction.}
The exact definitions are
\path{Formalism/LocalRingCardinality/DataTheta.lean},
$\operatorname{TPar}$ at 501--511, interlacing at 529--533, blocks at 535--536,
$D_P$ at 538--543, $X_P$ at 545--548 and centredness at 550--554.
The parity, gap and $\operatorname{Hyp3}$ definitions are in
\path{Formalism/SharpRestriction/CuspSpanNormBounds/CayleySignature.lean},
106--118.
The first package requires $t$ odd, $r+5\le t$, the indicated arrays and
$\eta_\tau\in\{1/2,-1/2\}$. Since $r'\le r$, its inequality weakens.
The interlacing conditions involving $u,v,b_0$ are unchanged; those involving
$b_j,e_j,b_{j+1}$ are merely restricted. Parity removes entries from integral
arrays. Every centredness condition retains its old block length $t$ and
either keeps its old indices or restricts them. Hence all these conditions
remain true.

Before its final translation, the source's signed set $D_P(j,\tau)$ is
\[
 \{(v_i,\mathrm{true})\}_{i<k-1}
 \ \cup\
 \{(-b_i,[\,i\le j\text{ or }i-j\text{ even}\,])\}_{i<M}
 \ \cup\
 \bigcup_{i<j}\{(x,\mathrm{true}):x\in\operatorname{block}(-e_i,t)\}.
\]
Natural subtraction in the Boolean test gives the identical answer on
each retained pair of indices. For $j\le M'$, the first set is unchanged,
the second loses precisely the indices $i\ge M'$, and the third is unchanged
because $i<j$ implies $i<M'$. Both records translate by the same $\eta_\tau$
when $j$ is odd. This proves the signed inclusion.
The set $X_P$ consists of the unchanged $u$-set and the unchanged blocks
of length $t+1$ with indices $i<j$, proving its asserted equality.

Every pair of opposite-sign entries in the smaller set therefore occurs in
the old set. Its old separation exceeds $(r+Mt)^3$, and
\[
 (r'+M't)^3\le(r+Mt)^3.
\]
Thus the strict gap inequality survives.
The remaining infinity-type identity transports automatically:
every source $\operatorname{ThetaX}$ satisfies $\operatorname{etaT\_pos}$
at \path{DataTheta.lean}, line 614, requiring infinity type $1/2$.
The old identity consequently says $\eta_\tau=-1/2$, which is also the
new package's required identity.
The only unaccounted $\operatorname{Hyp3}$ conjunct is
$1801\le r'+t$, precisely the displayed condition.
Since $r'=r-(M-M')$, the sufficient margin follows. This proves (a).
It does not identify the old indefinite $Y_{M'}$ with the rebuilt endpoint.

\textit{Finite nonzero-period induction.}
Here are the exact six inputs in \path{CayleySignature.lean}:
\path{RallisSiegelWeilFormula}, 154--161;
\path{LiftParameterConditions}, 163--166;
\path{ThetaUnfoldingFreeData}, 168--175;
\path{ThetaStablePeriodVanishing}, 177--184;
\path{RealThetaLiftDiscreteSeries}, 186--194; and
\path{RealConstantNonvanishing}, 196--199.
Their conjunction is \path{AdmS}, 201--202.
The conclusion \path{PointsLemma} is at 138--148 and the auxiliary
\path{PerNZ} at 150--152.
The name \path{TowerRallisLemma} at line 204 denotes the proposition of
the implication; it is not itself a theorem declaration.
We prove that implication from the six inputs here.

The first input supplies a character sequence $\mu$ and irreducible nonzero
representations $\pi_j$ of the definite groups, with their prescribed types.
For every $j<M$, every vector in $\pi_{j+1}$ is a finite sum of upward theta
lifts of vectors in $\pi_j$, using factorizable finite-type Schwartz data.
The second input supplies terminal $\phi_M\in\pi_M$ and a cuspidal vector
$F_M$ of the required type, with nonzero top period.

Suppose $\phi_{j+1},F_{j+1}$ have been constructed with these properties.
Expand $\phi_{j+1}$ as the finite sum supplied by the first input.
The period is additive in this vector; the conjugation in the lower
period is additive as well. Since the sum of period values is nonzero,
some summand has nonzero period. Choose its vector $\phi_j\in\pi_j$ and
upward Schwartz datum $\Phi'_j$.
The third input supplies a factorizable finite-type downward datum $\Phi_j$
whose lift will have nonzero \path{PerNZ}, once its cuspidality is supplied.

For that last requirement, let $W_{j+1}$ be the subspace of automorphic
vectors which are cuspidal and have the specified \path{IsDv} type.
It is a complex subspace: the latter condition says that the conjugate
vector lies in a fixed complex type submodule, and inverse images of
subspaces under a conjugate-linear additive map are subspaces.
It is stable under the finite and real group actions by the sixth input
and the cusp-stability fields of \path{ThetaX}.
For every vector in this subspace, the second clause of the fifth input
gives all lower-tower vanishings needed by the fourth input.
That fourth input therefore makes
\[
 F_j=v\theta_j(\Phi_j,F_{j+1})
\]
cuspidal. The first clause of the fifth input gives its \path{IsDv} type,
and the third input now gives its nonzero period against $\phi_j$.
This completes one step, hence the finite descent from $M$ to zero.
Indices outside the constrained finite range may be filled by zero.
The required character and $\pi_0$ came from the first input.

At zero, the terminal alternative in \path{PerNZ} is unavailable since
$M\ge2$. Thus a lower period is nonzero for some position and
Fourier--Jacobi data. Its rank difference is
\[
 \operatorname{rank}Y_0-\operatorname{rank}X_0
       =(k+M-1)-k=M-1,
\]
which is odd. The field \path{ThetaX.Bper_odd} at
\path{DataTheta.lean}, line 663, replaces the Fourier--Jacobi datum by zero
without changing the value. This is the exact zero-datum conclusion of
\path{PointsLemma}. The required additivity, stability and conjugation
identities are in the full \path{ThetaX} interface at 573--694.
No additional lower bound on the number of steps entered this proof.
This proves (b).

\textit{Actual endpoint period and its weaker field hypotheses.}
The source \path{Formalism/LocalToGlobal/Columns/Record.lean} takes arbitrary
rank, weight and endpoint-level inputs at 1434--1438, defines its row and
point data at 1440--1442 and proves point goodness at 1444--1448.
For a basis $\{h_a\}$ of the row invariant space and a point-class
representative $c$, coefficient pairing with $h_a(c)$ gives a linear
functional on the endpoint coefficient representation.
The evaluation map has linearly independent images of its basis.
The source's extension functional, at 154--205, consequently extends
this coefficient functional to the ambient complex function space and
agrees with it on evaluated vectors.

The definition \path{periodC}, 1568--1571, is the finite sum of these
functionals over point classes and row coordinates, with reciprocal
stabilizer-order weights.
The coefficient Gram matrix is invertible by
\path{Columns/RecordInst.lean}, \path{ptEv_gram_unit}, 47--58.
Thus each coefficient functional is represented by a vector under its
classwise inner product. The theorem
\path{Record.periodC_eq_ipCl}, 1580--1590, and its version in
\path{RecordInst.lean}, 60--64, identify the entire period with the
classwise inner product against \path{jMCL}.
All definitions use the general rank record, with endpoint dimension
$N_R(m)-1$.
The original \path{recPer}, 1603--1610, combines this map with evaluation
and the branching map \path{lowC}, 1558--1564, from $N_R(m)$ to $N_R(m)-1$.
It has no separate numerical existence assumption for a particular $m$.

More precisely, the argument $h_X:X.\operatorname{Ok}$ in the displayed
definition of \path{periodC} is used only through its point-goodness
proof \path{ptOk W \Lambda_M hX}.
Replace that proof by an explicit
\[
 h_{\rm pt}:(\operatorname{ptGSD}W\Lambda_M).\operatorname{Ok}
\]
and supply likewise
$h_{\rm row}:(\operatorname{rowGSD}W\Lambda_M).\operatorname{Ok}$.
Leave the same classwise functionals and row basis in place.
This defines a weakened endpoint map.
The same finite-sum proof gives its inner-product identity.
The Gram theorem just cited has no $X.\operatorname{Ok}$ argument,
so that input has not been reintroduced.

The two goodness proofs can actually be supplied under the weaker
conjugation assumptions in (c).
Read the entire proof \path{formGSD_ok} in
\path{Columns/ColFormsLemmas.lean}, 83--119.
Its five \path{GSD.Ok} fields are obtained as follows.
First, compatibility with complex conjugation at all embeddings of $E$
makes the standard hermitian form positive at all infinite places.
The source's decorated-lattice finiteness theorem supplies finitely many
representatives. The representative move is an isometry paired with the
identity level element; it already lies in \path{FormG}.
Thus this is a finite-orbit proof for the required group itself, not an
inference from finiteness for a larger group.
Second, positivity makes each decorated-lattice stabilizer finite.
The inclusion \path{FormG.incl} is injective, so its stabilizers are finite.
These are exactly the arguments at 84--101.

Third, every value of the finite level character is a root of unity,
by \path{LevFam.chr_root}, 77--81.
Such a scalar and its inverse are algebraic integers, so multiplication
by it preserves an integral lattice.
Combining this with the two supplied coefficient-lattice equivariance
identities proves the \path{nab_smul} and \path{del_smul} fields
(102--109).
Fourth, compatibility of one embedding $\sigma_0$ of the coefficient
field with conjugation gives $\bar c c=1$ for every such root of unity:
after the injective embedding its complex norm is one.
The invariant hermitian coefficient pairing remains invariant after
the character twist, giving the final \path{inv} field (110--119).
These steps cover all five fields; no quadratic-subfield hypothesis is used.

For point coefficients, \path{Record.ptOk}, 1444--1448, supplies the
needed coefficient equivariance from the round coefficient-lattice
identities and \path{roundPk_inv}; its only field input there is
\path{emb_star}.
For row coefficients,
\path{Columns/RowFormsLemmas.lean}, \path{rowVal_eqv}, 20--23,
uses the row lattice identities and \path{squarePk_inv} with precisely
the same input.
The asserted selected-embedding conjugation identity supplies it.
Hence both explicit goodness certificates exist.
For $E=K=\mathbf Q(\zeta_7)$, canonical complex conjugation commutes
with all Galois embeddings, so every one of these weaker hypotheses
holds for the selected embeddings $\sigma_1,\sigma_2,\sigma_3$.
This proves (c), for every specified admissible weight package and level.
The exact weight interfaces are \path{Columns/AbelianHostFieldConditions.lean},
\path{Wts} at 443--449, its goodness at 455--460 and endpoint weights
at 472--476.
The full original \path{recPer} is still attached to a
\path{recChain hX}; the preceding proof constructs the entire endpoint
period-map layer and does not silently supply the remaining global kernels.

\textit{The unchanged field wrapper.}
The carrier of \path{Fields.Kq} is displayed at
\path{AbelianHostFieldConditions.lean}, line 60, within its definition
at 59--78. The field \path{Fields.Ok.quad} at line 108 requires degree
two; \path{Fields.Core} still retains this requirement at line 157.
In $(\mathbf Z/7\mathbf Z)^\times$ the powers of $3$ are
$3,2,6,4,5,1$. Thus $\sigma_3$ generates the order-six Galois group.
Equality $\sigma_1(z)=\sigma_3(z)$ forces $z$ to be fixed by that group,
so $z\in\mathbf Q$; conversely every rational element is fixed by all
selected embeddings. An injective coefficient-field enlargement does not
change these equalities. Simultaneous conjugation or relabeling conjugates
the relative-automorphism subgroup, which remains the full group.
This proves the exact obstruction and its invariances.

For the general assertion, let $E/\mathbf Q$ be Galois CM of degree $2g$
and let the common field $L$ have degree two.
All $g$ selected embeddings restrict to the same embedding of $L$.
That embedding has exactly $[E:L]=g$ extensions, so the type consists
of all of them.
If $L$ were real, the $g$ complex-conjugate embeddings would have the same
restriction, giving $2g$ distinct extensions instead, a contradiction.
Hence $L$ is imaginary quadratic and the type is induced from it.
Conversely an induced type is a coset of $\operatorname{Gal}(E/L)$;
its relative automorphisms fill that subgroup, whose fixed field is $L$.

The contrasting type $\{1,2,4\}$ has an order-three relative subgroup.
Its fixed element $z=\zeta_7+\zeta_7^2+\zeta_7^4$ satisfies
$\bar z=-1-z$ and $z\bar z=2$, hence $z^2+z+2=0$.
Its common field is therefore $\mathbf Q(\sqrt{-7})$.
That is a different CM type, not a relabeling of the actual Jacobian type.
This proves (d).
\end{proof}

Two further source distinctions prevent the conclusion from overstating
what has been constructed.
First, \path{SpanPowerMembership/ThetaFamilyAdmissibleThree.lean}
does construct an abstract \path{thetaG} at 71--153, proves its six
statements at 215--274 and states \path{thetaG_exists} at 277--280.
Its spaces are a common auxiliary Hilbert space, its group actions are
identity and its periods are inner products.
This is a realization of the abstract interface, but does not identify
the actual mixed-CM theta germs with that realization's kernels.
Second, \path{SpanPowerMembership/NZColC.lean}, \path{periodN}, 160--167,
accepts an arbitrary bilinear form on the terminal class functions and
uses identity evaluation/lowering maps.
Its existence likewise supplies no geometric period identification.
The source's initial column nontriviality results
\path{Columns/NonDeg.lean}, \path{MramL_ne_bot}, 522--525, and
\path{finrank_MramV_pos_of_sizes}, 900--902, use $2fk<g_0-1$,
which stays unchanged when only $m$ decreases.
They are not nonvanishing statements for every terminal iterated image.

Together, DP44 and DP52 remove a rank, parameter, abstract induction and
endpoint-period existence obstruction to rebuilding a shorter even chain.
The remaining construction must identify one actual mixed-type theta
family with these periods, its algebraic theta germs and selected local
kernel operators; supply its six theta-correspondence identities; and
prove its local/global quantitative certificates.
The failed quadratic field wrapper is explicitly bypassed at the
endpoint layer, rather than treated as an intrinsic period obstruction.
The net deficit in DP44 must still be evaluated for that actual rebuilt
family. Neither ordinary primitive-path rates nor equal-direction averages
may be substituted for its selected gains, and no minimality of the old
lift count follows from this proposition.
